\documentclass{situationist-preprint}

\title{Consciousness Without Attention}
\msid{MND-24-a}
\msfamily{MND-24}
\msdiscipline{Philosophy of Mind}
\mstype{Principal manuscript}
\mskeywords{attention; phenomenal consciousness; inattentional blindness; visual periphery; overflow; introspection}

\begin{document}
\maketitle

\begin{abstract}
Is attention necessary for consciousness? Inattentional blindness and related phenomena suggest that what is not attended is not experienced, and theories on which consciousness is attended intermediate-level representation make the dependence constitutive. We argue that phenomenal consciousness extends beyond the scope of attention. The evidence for necessity is evidence about report, and report requires attention; so it does not discriminate between the hypothesis that unattended items are not experienced and the hypothesis that they are experienced but not reported. Three considerations then favor the second: the visual field is experienced as complete and without gaps, although attention covers a small part of it; partial-report experiments show that more is phenomenally present than can be attended and reported; and the gist and ensemble properties of unattended regions are perceived and used. We show, as a corollary, why the visual periphery escapes reliable introspection: introspective judgment of a state's character requires attention to it at the grain judged.
\end{abstract}

\section{The problem}

Subjects counting basketball passes often fail to notice a person in a gorilla suit crossing the scene \citep{SimonsChabris1999}; subjects attending to a task often fail to report an unexpected stimulus presented at fixation \citep{MackRock1998}. On one reading these results show that without attention there is no conscious experience. \citet{Prinz2012} built the reading into a theory: consciousness is the attended intermediate-level representation of the world. \citet{Cohen2012} argued that no experiment has shown conscious experience without attention. On the other side, \citet{Lamme2003} and \citet{KochTsuchiya2007} argued that attention and consciousness are distinct processes, and \citet{Block2007} argued that phenomenal consciousness overflows cognitive access. We argue that phenomenal consciousness extends beyond the scope of attention.

\section{Framework}

\begin{definition}[Attention; scope]\label{d:att}
\emph{Attention} is the selection of some perceptual information for enhanced processing and availability to report, reasoning and the control of action. The \emph{scope} of attention at a time is the information so selected.
\end{definition}

\begin{definition}[Phenomenal consciousness]\label{d:pc}
A state is \emph{phenomenally conscious} if there is something it is like for its subject to be in it.
\end{definition}

\begin{principle}[Report requires attention]\label{p:rep}
A subject can report, or make a fine-grained introspective judgment about, a state only if information about it is within the scope of attention at the grain reported.
\end{principle}

\begin{principle}[Completeness of the field]\label{p:field}
In normal waking vision, the visual field is phenomenally present as a whole: every region of it has some phenomenal character at each time, and no region is experienced as a gap.
\end{principle}

Principle~\ref{p:field} does not say that every region is experienced in detail. It says that the periphery is not phenomenally absent: one does not experience a spotlight surrounded by nothing, and the regions outside attention are experienced as colored, textured and populated at a coarse grain. The principle is a claim about the structure of visual experience that the alternative hypothesis must explain away, not a report of introspected detail.

\section{The argument}

\begin{lemma}\label{l:under}
The evidence from inattentional blindness does not discriminate between the hypotheses that unattended items are not experienced and that they are experienced but not reported.
\end{lemma}

\begin{proof}
The evidence consists of subjects' failures to report unattended items. By Principle~\ref{p:rep}, unattended items are not reportable at the grain required. So both hypotheses predict the failure to report. Inattentional blindness experiments also probe memory after the event, and the second hypothesis predicts failure there too, since items not attended are not encoded for later report \citep{Wolfe1999}.
\end{proof}

\begin{theorem}\label{t:main}
Phenomenal consciousness extends beyond the scope of attention.
\end{theorem}

\begin{proof}
By Lemma~\ref{l:under}, the evidence for the necessity of attention is neutral. Three considerations favor the extension of phenomenal consciousness beyond attention. First, by Principle~\ref{p:field}, every region of the visual field has phenomenal character at each time, while attention, by the standard estimates, selects only a few items or a small region at once. So some regions with phenomenal character lie outside its scope. Second, in partial-report experiments, subjects shown a brief array of letters can report only a few, but can report any row cued after the array has disappeared \citep{Sperling1960}; they report having seen all the letters, and their performance shows that information about all of them was available for a short time, before attention selected some for report. The best explanation is that the whole array was phenomenally present and attention selected from what was present. Third, subjects extract the gist of natural scenes and the average properties of ensembles in regions to which they do not attend, and can categorize scenes when attention is engaged by a demanding central task \citep{Li2002}; these properties are part of the coarse character the unattended periphery is experienced as having. An attention-necessity theory must say that the peripheral field is phenomenally empty, contrary to Principle~\ref{p:field}, or that it is attended, contrary to the limited capacity of attention.
\end{proof}

\begin{corollary}\label{c:periph}
Introspection of the unattended periphery is unreliable at any grain finer than the coarse grain at which it is attended, and becomes reliable when attention is directed to it.
\end{corollary}

\begin{proof}
By Theorem~\ref{t:main}, there are phenomenal states at the periphery outside the scope of attention. By Principle~\ref{p:rep}, fine-grained introspective judgments about them cannot be formed while they remain unattended; a judgment about their detail can only be reached by inference from expectation or theory, which is not the method on which introspective reliability rests. When attention is directed to a peripheral point, the state comes within its scope and a judgment fixed by the state itself becomes possible.
\end{proof}

Corollary~\ref{c:periph} is the explanation that MND-20 invokes for Schwitzgebel's observation that people are often wrong about how clear the visual field is away from fixation \citep{Schwitzgebel2008}. The errors concern states outside attention; they do not show that introspection is unreliable about attended current states.

\section{Objections and replies}

\begin{objection}[Refrigerator light]\label{o:fridge}
Whenever we check whether the periphery is experienced, we attend to it, so the sense of a complete field may be an illusion produced by the availability of attention: like the refrigerator light, it is always on when we look.
\end{objection}

\begin{reply}
The objection grants Principle~\ref{p:field} as a datum and explains it away. But the explanation predicts that unattended regions are experienced as nothing at all, and nothing in visual experience, even under heavy central load, is like a gap; the subject who misses the gorilla does not experience a hole where it walked. The partial-report evidence also concerns information available before attention selects it, which the refrigerator hypothesis does not explain.
\end{reply}

\begin{objection}[Diffuse attention]\label{o:diffuse}
The periphery is attended, diffusely; so there is no consciousness without attention \citep{Prinz2012}.
\end{objection}

\begin{reply}
If diffuse attention is selection for enhanced processing, its capacity limits apply and it cannot cover the whole field at once. If it is not, it is not attention in the sense of Definition~\ref{d:att}, and the claim that consciousness requires it becomes the claim that consciousness requires being processed at all, which no one disputes.
\end{reply}

\begin{objection}[Unfalsifiability]\label{o:falsi}
If unattended experience cannot be reported, the hypothesis that it exists is beyond empirical test.
\end{objection}

\begin{reply}
The hypothesis is tested indirectly, by its explanatory power with respect to the partial-report data, gist perception and the structure of the visual field, as other hypotheses about unobservables are tested.
\end{reply}

\section{Conclusion}

The evidence that attention is necessary for consciousness is evidence that it is necessary for report. Once this is recognized, the structure of the visual field, the partial-report data and the perception of gist favor the conclusion that phenomenal consciousness extends beyond attention, and the same conclusion explains why the periphery escapes reliable introspection.

\begin{thebibliography}{99}\small
\bibitem[Block(2007)]{Block2007} Block, N. (2007). Consciousness, accessibility, and the mesh between psychology and neuroscience. \emph{Behavioral and Brain Sciences}, 30(5--6), 481--499.
\bibitem[Cohen et~al.(2012)]{Cohen2012} Cohen, M.~A., Cavanagh, P., Chun, M.~M. and Nakayama, K. (2012). The attentional requirements of consciousness. \emph{Trends in Cognitive Sciences}, 16(8), 411--417.
\bibitem[Koch and Tsuchiya(2007)]{KochTsuchiya2007} Koch, C. and Tsuchiya, N. (2007). Attention and consciousness: Two distinct brain processes. \emph{Trends in Cognitive Sciences}, 11(1), 16--22.
\bibitem[Lamme(2003)]{Lamme2003} Lamme, V.~A.~F. (2003). Why visual attention and awareness are different. \emph{Trends in Cognitive Sciences}, 7(1), 12--18.
\bibitem[Li et~al.(2002)]{Li2002} Li, F.~F., VanRullen, R., Koch, C. and Perona, P. (2002). Rapid natural scene categorization in the near absence of attention. \emph{Proceedings of the National Academy of Sciences}, 99(14), 9596--9601.
\bibitem[Mack and Rock(1998)]{MackRock1998} Mack, A. and Rock, I. (1998). \emph{Inattentional Blindness}. MIT Press.
\bibitem[Prinz(2012)]{Prinz2012} Prinz, J.~J. (2012). \emph{The Conscious Brain: How Attention Engenders Experience}. Oxford University Press.
\bibitem[Schwitzgebel(2008)]{Schwitzgebel2008} Schwitzgebel, E. (2008). The unreliability of naive introspection. \emph{The Philosophical Review}, 117(2), 245--273.
\bibitem[Simons and Chabris(1999)]{SimonsChabris1999} Simons, D.~J. and Chabris, C.~F. (1999). Gorillas in our midst: Sustained inattentional blindness for dynamic events. \emph{Perception}, 28(9), 1059--1074.
\bibitem[Sperling(1960)]{Sperling1960} Sperling, G. (1960). The information available in brief visual presentations. \emph{Psychological Monographs: General and Applied}, 74(11), 1--29.
\bibitem[Wolfe(1999)]{Wolfe1999} Wolfe, J.~M. (1999). Inattentional amnesia. In V. Coltheart (ed.), \emph{Fleeting Memories: Cognition of Brief Visual Stimuli}. MIT Press.
\end{thebibliography}
\end{document}
